\begin{table}[!tbp]
\centering
\caption{PACE~2021 exact track, the 173 instances with known optimum: lower
bounds.  Checked (values returned by the checker from archived
certificates): the ruin-and-recreate star packing, the CertiFlip bound
(budget \SI{60}{s}; measured times in Section~\ref{sec:exp-time}), the
CertiFlip bound with \SI{1800}{s} (every fourth instance), the star local
search (\SI{60}{s}) and the larger of the CertiFlip bound and the star local
search.  Not certified: the triangle
packings, computed earlier on a 4-core workstation; the metric LP on $P$, a
floating-point LP value reported only where cutting planes converged within
\SI{900}{s} (every fourth instance); the two B\&B rows, the root bounds of the
KaPoCE branch-and-bound as published by its authors (reproduced by a rerun
without data reductions, see the text); and the last row, which
takes the larger of a checked and a published value.}
\label{tab:pace-bounds}
\footnotesize
\begin{tabular}{lrrrr}
\toprule
bound & instances & mean LB/OPT & min LB/OPT & LB $=$ OPT \\
\midrule
greedy triangle packing & 173 & 0.8612 & 0.660 & 14 \\
fractional triangle packing & 173 & 0.9664 & 0.811 & 30 \\
ruin-and-recreate star packing & 173 & 0.8712 & 0.247 & 17 \\
CertiFlip bound & 173 & 0.9874 & 0.816 & 106 \\
CertiFlip bound, 1800 s & 45 & 0.9948 & 0.940 & 34 \\
star local search & 173 & 0.9984 & 0.960 & 122 \\
metric LP on $P$, triangle rows & 39 & 0.9762 & 0.814 & 21 \\
B\&B star packing~\cite{BlasiusEtAl22sea} & 173 & 0.9948 & 0.891 & 79 \\
B\&B $P_3$ packing~\cite{BlasiusEtAl22sea} & 173 & 0.9742 & 0.813 & 59 \\
larger of CertiFlip and star local search & 173 & 0.9989 & 0.972 & 130 \\
larger of both and B\&B star & 173 & 0.9989 & 0.972 & 131 \\
\bottomrule
\end{tabular}

\end{table}

\begin{table}[!tbp]
\centering
\caption{Mean LB/OPT on the solved PACE exact instances by edge density
$m/\binom n2$.  The fractional triangle packing was computed earlier on a
4-core workstation.}
\label{tab:pace-density}
\footnotesize
\setlength{\tabcolsep}{4pt}
\begin{tabular}{lrrrrr}
\toprule
edge density & instances & CertiFlip & sparse star & fractional & B\&B star \\
 & & bound & packing & triangle packing & packing \\
\midrule
$[0,0.1)$ & 35 & 0.9978 & 0.9529 & 0.9614 & 0.9934 \\
$[0.1,0.3)$ & 40 & 0.9781 & 0.8623 & 0.9646 & 0.9971 \\
$[0.3,0.5)$ & 48 & 0.9823 & 0.9046 & 0.9448 & 0.9914 \\
$[0.5,1)$ & 50 & 0.9923 & 0.7892 & 0.9922 & 0.9974 \\
\bottomrule
\end{tabular}

\end{table}

\paragraph{Strength of the bounds.}  The star local search, run for
\SI{60}{s} per instance, gives the strongest bound on this track
(Table~\ref{tab:pace-bounds}).  Its mean LB/OPT on the 173 solved instances
is \numPaceSlsMean{}, against \numPaceMeanStar{} for the root star packing of
the KaPoCE branch-and-bound as published and \numPaceMeanOurs{} for the
CertiFlip bound, and it is at least \numPaceSlsMin{} of the optimum on every
instance.  It proves optimality on \numPaceSlsOpt{} instances, against 79 for
the B\&B star packing, 59 for its $P_3$ packing and \numPaceOptOurs{} for the
CertiFlip bound, and together with the CertiFlip bound, both checked, on
\numPaceBestOpt{} (mean \numPaceBestMean{}).  Against the published B\&B star
packing it is larger on \numPaceSlsAbove{} instances, equal on
\numPaceSlsEqual{} and smaller on \numPaceSlsBelow{}; adding that bound to
the two checked ones proves optimality on only one more instance
(\numPaceBestMaxOpt{}).  The advantage holds in every density band
(Table~\ref{tab:pace-density}).  The CertiFlip bound is larger than the star
local search on \numPaceSlsUnderCf{} instances, by one to eight units, and
proves optimality on eight of them, where the star local search does not:
there the block LPs close the last units.  Below density $0.1$ the two are
equal on average.  Each run took about the budget (median \numPaceSlsTimeMedian{}\,s,
at most \numPaceSlsTimeMax{}\,s, including the certificate).

All three star packings pack the same induced stars, and our earlier
ruin-and-recreate packing reaches only \numPacePackMean{} on average
(Table~\ref{tab:pace-bounds}): the difference lies in the packing
heuristic.  The published data state neither the time
spent on the root bounds nor whether they were computed after the data
reductions of the branch-and-bound, so we reran its lower-bound program on
the input graphs, without reductions (the program of
Section~\ref{sec:exp-cert}, \SI{900}{s} limit, one run per instance;
\numKrootPaceColab{} on the Colab machines, one at a time, the others on a
4-core cloud machine with a \SI{2.2}{GHz} Xeon, at most three at a time).
The star packing finished within the limit on \numKrootPaceStarN{} of the
\numKrootPaceN{} instances (\numKrootPaceTimeoutList{} stopped after the $P_3$
packing).  Its $P_3$ packing equals the published value on
\numKrootPacePthreeEqual{} of \numKrootPacePthreeN{}, its star packing on
\numKrootPaceStarEqual{}, and on the others the star packing differs by at
most \numKrootPaceStarDiffMax{} (\numKrootPaceStarAbove{} larger,
\numKrootPaceStarBelow{} smaller, at most \numKrootPaceStarRelMax{}; the
heuristic uses a random generator, which we ran with its default seed, and
the published values may come from another version).  The star bound took a
median of \numKrootPaceTimeMedian{}\,s and at most \numKrootPaceTimeMax{}\,s.
The published root bounds are therefore bounds of the input graphs.  The
heuristic stops by its own rule, after five times as many rounds as improved
the packing; the star local search instead runs until its time limit.

On every fourth instance we also ran the bound with \SI{1800}{s} and one
block covering the graph, and the metric LP on $P$ with triangle rows only.
With the longer budget the mean LB/OPT on these \numPaceWarmN{} instances is
\numPaceWarmMean{}, against \numPaceWarmOursMean{} with the \SI{60}{s} runs
and \numPaceWarmStarMean{} for the B\&B star packing, and the bound is
optimal on \numPaceWarmOpt{} of them (\numPaceWarmOursOpt{} at \SI{60}{s}).
The triangle LP converged within \SI{900}{s} on \numPaceTriN{} of these
instances, with a mean LB/OPT of \numPaceTriMean{} (\numPaceTriOursMean{}
for our bound on the same instances).  Of the \numWarmRuns{} bound runs,
\numWarmLocal{} were made on a laptop (\numCpuLocal), several at a time, and one
on a Colab machine.  The limit applies to the LP phase; the star-packing warm
start before it (sized to about 30\% of the limit) and the certificate come
on top, and \numWarmOver{} of the \numWarmRuns{} runs took longer than
\SI{1800}{s}, the longest \numWarmMax{}\,s.  The checker accepted all their
certificates.

\paragraph{Open instances.}  Of the \numPaceOpen{} instances of the exact
track without a known optimum, our best checked bound exceeds the larger
published root bound~\cite{BlasiusEtAl22sea} on \numPaceOpenImproved{}
(Table~\ref{tab:pace-open}); on \texttt{exact179} it leaves a gap of one to
the best known solution.  On the other \numPaceOpenBelow{} it is smaller than
the published root bound, by at most \numPaceOpenMaxBelow{}.  None of the
open instances is closed.

\begin{table}[!tbp]
\centering
\caption{Open instances of the PACE~2021 exact track on which the checked
lower bound improves the published one.  Published bounds: the best upper
bound and the larger of the two root bounds of~\cite{BlasiusEtAl22sea}; our
UB: the CertiFlip run of Table~\ref{tab:pace-bounds}; checked LB: the largest
checked bound of our runs.}
\label{tab:pace-open}
\footnotesize
\begin{tabular}{lrrrrrr}
\toprule
instance & $n$ & $m$ & published UB & published LB & our UB & checked LB \\
\midrule
\texttt{exact042} & 100 & 1563 & \num{1055} & \num{1006} & \num{1055} & \num{1017} \\
\texttt{exact043} & 100 & 1624 & \num{984} & \num{968} & \num{984} & \num{978} \\
\texttt{exact051} & 120 & 2251 & \num{1619} & \num{1508} & \num{1619} & \num{1533} \\
\texttt{exact052} & 120 & 2448 & \num{1473} & \num{1431} & \num{1473} & \num{1442} \\
\texttt{exact060} & 124 & 2288 & \num{1491} & \num{1427} & \num{1491} & \num{1446} \\
\texttt{exact069} & 140 & 2915 & \num{2140} & \num{1994} & \num{2140} & \num{2015} \\
\texttt{exact070} & 140 & 3014 & \num{1792} & \num{1762} & \num{1792} & \num{1779} \\
\texttt{exact082} & 160 & 4015 & \num{2434} & \num{2397} & \num{2434} & \num{2414} \\
\texttt{exact083} & 160 & 4259 & \num{2939} & \num{2808} & \num{2941} & \num{2824} \\
\texttt{exact091} & 180 & 4602 & \num{3234} & \num{3072} & \num{3236} & \num{3087} \\
\texttt{exact092} & 180 & 6067 & \num{3563} & \num{3520} & \num{3563} & \num{3529} \\
\texttt{exact101} & 200 & 6731 & \num{4198} & \num{4135} & \num{4198} & \num{4148} \\
\texttt{exact102} & 200 & 6746 & \num{3925} & \num{3895} & \num{3925} & \num{3903} \\
\texttt{exact103} & 200 & 7081 & \num{4572} & \num{4381} & \num{4572} & \num{4382} \\
\texttt{exact105} & 200 & 12309 & \num{5320} & \num{5314} & \num{5320} & \num{5318} \\
\texttt{exact179} & 330 & 1370 & \num{633} & \num{620} & \num{633} & \num{632} \\
\texttt{exact180} & 330 & 2256 & \num{1081} & \num{1052} & \num{1083} & \num{1073} \\
\texttt{exact181} & 330 & 3253 & \num{1563} & \num{1513} & \num{1568} & \num{1532} \\
\texttt{exact182} & 330 & 4329 & \num{2159} & \num{2091} & \num{2165} & \num{2100} \\
\texttt{exact183} & 330 & 5613 & \num{2789} & \num{2717} & \num{2788} & \num{2735} \\
\texttt{exact184} & 330 & 7270 & \num{3684} & \num{3582} & \num{3692} & \num{3597} \\
\texttt{exact195} & 620 & 3992 & \num{2529} & \num{2450} & \num{2533} & \num{2476} \\
\bottomrule
\end{tabular}

\end{table}

\begin{table}[!tbp]
\centering
\caption{PACE~2021 exact track (200 instances): primal quality against the
best known value (published upper bound or better value of any run).
CertiFlip: the certificate runs of Table~\ref{tab:pace-bounds}, one seed on
the Colab machines; SCC: one seed on the complete encoding, on a 4-core
machine with three runs in parallel; $^a$: best of three seeds from earlier
runs on the same 4-core machine with four runs in parallel.}
\label{tab:pace-primal}
\footnotesize
\begin{tabular}{lrrr}
\toprule
method & best known & mean gap (\%) & max gap (\%) \\
\midrule
CertiFlip (60 s) & 192/200 & 0.006 & 0.32 \\
SCC~\cite{hausberger2025scc} (60 s) & 200/200 & 0.000 & 0.00 \\
KaPoCE (60 s)$^a$ & 200/200 & 0.000 & 0.00 \\
IteratedFlip$^a$ & 187/200 & 0.007 & 0.32 \\
Leiden-CPM$^a$ & 158/200 & 0.119 & 3.00 \\
Pivot$^a$ & 13/200 & 32.635 & 229.54 \\
\bottomrule
\end{tabular}

\end{table}

\paragraph{Primal quality.}  On the exact track the primal side is not the
bottleneck (Table~\ref{tab:pace-primal}).  CertiFlip finds the best known
value on \numCfHit{} of the 200 instances although it spends a large part
of its \SI{60}{s} on the bound; SCC on the complete encoding and KaPoCE find it
on all of them.  The certified gaps on this track are therefore determined by
the bound.
